\documentclass[11pt]{article}
\usepackage[margin=1in]{geometry}
\usepackage{amsmath,amssymb,amsthm,mathtools}
\usepackage{microtype}
\usepackage{booktabs}
\usepackage[hidelinks]{hyperref}
\usepackage{enumitem}

\newtheorem{theorem}{Theorem}[section]
\newtheorem{proposition}[theorem]{Proposition}
\newtheorem{corollary}[theorem]{Corollary}
\newtheorem{lemma}[theorem]{Lemma}
\theoremstyle{remark}
\newtheorem{remark}[theorem]{Remark}

\newcommand{\R}{\mathbb{R}}
\newcommand{\C}{\mathcal{C}}
\newcommand{\ip}[2]{\left\langle #1,#2\right\rangle_\eta}

\title{\textbf{Casimir Completion and Primitive Null Diamonds}\\
\large Edge centering, cyclotomic normalization, and the ramified Boolean return}
\author{Jeremy Ebert}
\date{2026}

\begin{document}
\maketitle

\begin{abstract}
The quadratic Casimirs of the rank-one ladder algebras $SU(2)$ and $SU(1,1)$ admit the familiar completions
\[
j(j+1)+\frac14=\left(j+\frac12\right)^2,
\qquad
k(k-1)+\frac14=\left(k-\frac12\right)^2.
\]
This note gives the common $\tfrac14$ a concrete Lorentz-cone realization.  On the factor cone
\[
T^2-X^2-Y^2=0,
\qquad
u=T+X,\quad v=T-X,\quad Y^2=uv,
\]
a primitive unimodular pair of spinors determines two null vectors whose midpoint has Lorentz norm squared $1/4$.  At general lattice spacing $\delta$, the midpoint norm is $\delta^2/4$.  The exact $SU(2)$ and positive-discrete-series $SU(1,1)$ ladder amplitudes can then be written in the unified edge-centered form
\[
A_\pm^2
=
\sigma\left(C+M_F^2-Z_\pm^2\right),
\qquad
M_F^2=\frac{\delta^2}{4},
\]
with $\sigma=+1$ for $SU(2)$ and $\sigma=-1$ for $SU(1,1)$.  Equivalently, the edge-centered representation radius satisfies
\[
R^2=C+M_F^2.
\]
The result does not assert that Farey arithmetic causes the Weyl or Casimir shift.  Rather, it identifies the shift with the same centered-null-edge quadratic invariant that appears canonically in the arithmetic factor-cone geometry.  A four-corner oscillator transition cell makes the half-step origin of the invariant explicit.  We then compare this cone quarter with the distinct cyclotomic normalization
\[
\mathcal Z=\frac{H+iI}{2},\qquad H^2=3I,
\]
for which $\mathcal Z\overline{\mathcal Z}=I$, and record the mod-$2$ return of the associated discriminant-$12$ lattice as factor exchange on a four-state Boolean cell.  These are exact constructions sharing a dyadic mechanism type; they are not asserted to be the same operator, norm, or lattice.
\end{abstract}

\section{The factor cone}

Let
\[
u,v>0,\qquad
X=\frac{u-v}{2},\qquad
T=\frac{u+v}{2},\qquad
Y=\sqrt{uv}.
\]
Then
\begin{equation}
\boxed{T^2-X^2-Y^2=0.}
\label{eq:cone}
\end{equation}
We use the Lorentz bilinear form
\[
\ip{(X,Y,T)}{(X',Y',T')}
=
TT'-XX'-YY'.
\]
Thus \eqref{eq:cone} is the future null cone.  The inverse light-cone coordinates are
\[
u=T+X,\qquad v=T-X,
\]
and the null condition becomes
\[
Y^2=uv.
\]

A convenient square-root parametrization is obtained from a real spinor
\[
\psi=
\begin{pmatrix}\alpha\\ \beta\end{pmatrix}.
\]
Define
\begin{equation}
Q(\alpha,\beta)
=
\left(
\frac{\alpha^2-\beta^2}{2},
\alpha\beta,
\frac{\alpha^2+\beta^2}{2}
\right).
\label{eq:spinormap}
\end{equation}
Then $Q(\psi)^2=0$, while
\[
u=\alpha^2,\qquad v=\beta^2.
\]
Equivalently,
\[
\psi\psi^T
=
\begin{pmatrix}
T+X&Y\\
Y&T-X
\end{pmatrix},
\]
whose determinant is precisely $T^2-X^2-Y^2$.

\section{Primitive unimodular null edges}

The arithmetic ingredient is elementary but rigid.  Take two spinors
\[
\psi_1=\begin{pmatrix}a\\b\end{pmatrix},
\qquad
\psi_2=\begin{pmatrix}c\\d\end{pmatrix},
\]
and let
\[
\Delta=\det(\psi_1,\psi_2)=ad-bc.
\]
Write $Q_i=Q(\psi_i)$.

\begin{lemma}[Determinant--Lorentz identity]
\label{lem:detlor}
For the quadratic map \eqref{eq:spinormap},
\[
\boxed{
\ip{Q_1}{Q_2}
=
\frac12(ad-bc)^2
=
\frac{\Delta^2}{2}.
}
\]
\end{lemma}

\begin{proof}
Direct expansion gives
\begin{align*}
\ip{Q_1}{Q_2}
&=
\frac{(a^2+b^2)(c^2+d^2)}{4}
-\frac{(a^2-b^2)(c^2-d^2)}{4}
-abcd\\
&=
\frac12(a^2d^2+b^2c^2-2abcd)
=
\frac12(ad-bc)^2.
\end{align*}
\end{proof}

When $\psi_1,\psi_2$ are primitive integer spinors with $|\Delta|=1$, and their projective ratios are written in standard reduced representatives with a common denominator orientation, those ratios are Farey adjacent.  We therefore call the corresponding pair $(Q_1,Q_2)$ a \emph{primitive unimodular null edge}.

Define its midpoint and half-difference by
\[
M_F=\frac{Q_1+Q_2}{2},
\qquad
D_F=\frac{Q_1-Q_2}{2}.
\]

\begin{proposition}[Null-edge midpoint invariant]
\label{prop:midpoint}
For any nondegenerate pair,
\[
M_F^2=\frac{\Delta^2}{4},
\qquad
D_F^2=-\frac{\Delta^2}{4},
\qquad
\ip{M_F}{D_F}=0.
\]
In particular, for a primitive unimodular edge,
\begin{equation}
\boxed{
M_F^2=\frac14.
}
\label{eq:quarter}
\end{equation}
If all cone coordinates are scaled by a lattice spacing $\delta$, then
\begin{equation}
\boxed{
M_F^2=\frac{\delta^2}{4}.
}
\label{eq:deltaquarter}
\end{equation}
\end{proposition}

\begin{proof}
Because $Q_1^2=Q_2^2=0$,
\[
M_F^2
=
\frac14(Q_1+Q_2)^2
=
\frac12\ip{Q_1}{Q_2}
=
\frac{\Delta^2}{4}.
\]
The formulas for $D_F^2$ and $\ip{M_F}{D_F}$ follow identically.  Scaling the vectors by $\delta$ multiplies all squared norms by $\delta^2$.
\end{proof}

Equation \eqref{eq:deltaquarter} is the geometric quantity to be compared with the rank-one Casimir completions below.

\section{Edge-centered ladder coordinates}

Let adjacent weights be separated by a spacing $\delta>0$.  For a transition leaving a weight $q$, define the edge-centered coordinate
\begin{equation}
\boxed{
Z_\pm=q\pm\frac{\delta}{2}.
}
\label{eq:Zedge}
\end{equation}
The elementary identity
\begin{equation}
q(q\pm\delta)
=
\left(q\pm\frac{\delta}{2}\right)^2-\frac{\delta^2}{4}
\label{eq:complete}
\end{equation}
is the algebraic source of the half-cell term.  Proposition~\ref{prop:midpoint} supplies a canonical arithmetic Lorentz realization of the same quantity:
\[
\frac{\delta^2}{4}=M_F^2.
\]

The point of the construction is not merely that two formulas contain the same number.  Both are edge-centering operations in the same quadratic cone geometry: \eqref{eq:complete} centers adjacent weights, while \eqref{eq:deltaquarter} centers adjacent primitive null rays.

\section{$SU(2)$: compact completion}

For an $SU(2)$ representation with highest weight $J$ and spacing $\delta$, write the quadratic Casimir eigenvalue as
\begin{equation}
C_+=J(J+\delta).
\label{eq:su2cas}
\end{equation}
At unit spacing $J=j$ and $C_+=j(j+1)$.

Define the edge-centered representation radius
\begin{equation}
R_+=J+\frac{\delta}{2}.
\label{eq:Rplus}
\end{equation}
Then
\begin{equation}
R_+^2
=
C_++\frac{\delta^2}{4}.
\label{eq:su2complete}
\end{equation}

\begin{theorem}[$SU(2)$ Casimir--null-edge completion]
\label{thm:su2}
Let $M_F$ be the midpoint of a scaled primitive unimodular null edge of spacing $\delta$.  Then
\begin{equation}
\boxed{
R_+^2=C_++M_F^2.
}
\label{eq:su2bridge}
\end{equation}
For a raising transition from weight $q$,
\begin{equation}
A_+^2
=
(J-q)(J+q+\delta)
=
R_+^2-Z_+^2,
\label{eq:su2amp}
\end{equation}
where $Z_+=q+\delta/2$.  Hence the transition is represented by the null cone point
\begin{equation}
\boxed{
Q_{J+}
=
(X,Y,T)
=
\left(
Z_+,\,
A_+,\,
\sqrt{C_++M_F^2}
\right).
}
\label{eq:su2point}
\end{equation}
\end{theorem}

\begin{proof}
Equation \eqref{eq:su2bridge} follows from \eqref{eq:su2complete} and $M_F^2=\delta^2/4$.  For the amplitude,
\begin{align*}
A_+^2
&=
(J-q)(J+q+\delta)\\
&=
J(J+\delta)-q(q+\delta)\\
&=
C_+-\left(Z_+^2-\frac{\delta^2}{4}\right)\\
&=
C_++M_F^2-Z_+^2\\
&=
R_+^2-Z_+^2.
\end{align*}
Therefore $T^2-X^2-Y^2=R_+^2-Z_+^2-A_+^2=0$.
\end{proof}

For lowering, $Z_-=q-\delta/2$ and the same calculation gives
\[
A_-^2=R_+^2-Z_-^2.
\]
At unit spacing this is the standard exact coefficient
\[
A_+^2=(j-m)(j+m+1)
=
\left(j+\frac12\right)^2
-
\left(m+\frac12\right)^2.
\]

Thus the compact ladder lies on the fixed-$T$ circle
\[
X^2+Y^2=R_+^2.
\]

\section{$SU(1,1)$: noncompact completion}

For the positive discrete series of $SU(1,1)$, let $K$ denote the Bargmann-index scale with spacing $\delta$ and write
\begin{equation}
C_-=K(K-\delta).
\label{eq:su11cas}
\end{equation}
At unit spacing $K=k$ and $C_-=k(k-1)$.

Define
\begin{equation}
R_-=K-\frac{\delta}{2}.
\label{eq:Rminus}
\end{equation}
Then
\begin{equation}
R_-^2=C_-+\frac{\delta^2}{4}.
\label{eq:su11complete}
\end{equation}

\begin{theorem}[$SU(1,1)$ Casimir--null-edge completion]
\label{thm:su11}
With the same scaled primitive null-edge midpoint $M_F$,
\begin{equation}
\boxed{
R_-^2=C_-+M_F^2.
}
\label{eq:su11bridge}
\end{equation}
For a raising transition from a $K_0$ weight $q$,
\begin{equation}
A_+^2
=
q(q+\delta)-C_-
=
Z_+^2-R_-^2,
\label{eq:su11amp}
\end{equation}
where $Z_+=q+\delta/2$.  Hence
\begin{equation}
\boxed{
Q_{K+}
=
(X,Y,T)
=
\left(
\sqrt{C_-+M_F^2},\,
A_+,\,
Z_+
\right)
}
\label{eq:su11point}
\end{equation}
is null.
\end{theorem}

\begin{proof}
Equation \eqref{eq:su11bridge} follows from \eqref{eq:su11complete} and Proposition~\ref{prop:midpoint}.  Using \eqref{eq:complete},
\begin{align*}
A_+^2
&=
q(q+\delta)-C_-\\
&=
Z_+^2-\frac{\delta^2}{4}-C_-\\
&=
Z_+^2-(C_-+M_F^2)\\
&=
Z_+^2-R_-^2.
\end{align*}
Therefore
\[
T^2-X^2-Y^2
=
Z_+^2-R_-^2-A_+^2=0.
\]
\end{proof}

At unit spacing, write $q=k+n$.  Then
\[
A_+^2
=
(n+1)(2k+n)
=
\left(k+n+\frac12\right)^2
-
\left(k-\frac12\right)^2.
\]
The noncompact ladder therefore lies on the fixed-$X$ hyperbola
\[
T^2-Y^2=R_-^2.
\]

\section{Unified Casimir--null-edge theorem}

The compact and noncompact cases differ only in which cone coordinate carries the completed Casimir radius.

\begin{theorem}[Casimir--null-edge completion]
\label{thm:unified}
Let $\delta>0$ be the weight spacing and let $M_F$ be the midpoint of a scaled primitive unimodular null edge, so that
\[
M_F^2=\frac{\delta^2}{4}.
\]
For the rank-one ladder systems above, define the edge-centered representation radius by
\[
R^2=C+M_F^2.
\]
Then the exact ladder amplitudes obey
\begin{equation}
\boxed{
A_\pm^2
=
\sigma\left(R^2-Z_\pm^2\right)
=
\sigma\left(C+M_F^2-Z_\pm^2\right),
}
\label{eq:unified}
\end{equation}
where
\[
Z_\pm=q\pm\frac{\delta}{2},
\qquad
\sigma=
\begin{cases}
+1,&SU(2),\\
-1,&SU(1,1).
\end{cases}
\]
Equivalently,
\begin{equation}
\boxed{
C=R^2-M_F^2.
}
\label{eq:deficit}
\end{equation}
\end{theorem}

Equation \eqref{eq:deficit} is the most compact geometric statement: the quadratic Casimir is the edge-centered radius squared minus one universal half-edge square.

The cone embedding is summarized by
\[
\begin{array}{c|ccc}
&X&Y&T\\
\hline
SU(2)&Z&A&R\\
SU(1,1)&R&A&Z
\end{array}
\]
with $T^2-X^2-Y^2=0$ in both cases.  The compact/noncompact change is therefore a permutation of the roles of $T$ and $X$:
\[
SU(2):\quad X^2+Y^2=R^2,
\]
\[
SU(1,1):\quad T^2-Y^2=R^2.
\]

\section{The four-corner oscillator transition cell}

The half-edge interpretation becomes particularly concrete in the two-oscillator realization.  Let
\[
|p,q\rangle,\qquad p,q\in\mathbb{Z}_{\ge0},
\]
be a two-mode number state, and attach to it the factor cell
\[
[\delta p,\delta(p+1)]
\times
[\delta q,\delta(q+1)].
\]
Its center is
\[
(u_c,v_c)
=
\delta\left(p+\frac12,q+\frac12\right).
\]

The four standard quadratic ladder amplitudes occupy the four vertices:
\[
\begin{array}{c|c|c}
\text{operator}&(u,v)&A^2\\
\hline
K_-&\delta(p,q)&\delta^2pq\\
J_+&\delta(p+1,q)&\delta^2(p+1)q\\
J_-&\delta(p,q+1)&\delta^2p(q+1)\\
K_+&\delta(p+1,q+1)&\delta^2(p+1)(q+1).
\end{array}
\]

Under
\[
X=\frac{u-v}{2},
\qquad
T=\frac{u+v}{2},
\]
the cell center is
\[
X_c=\frac{\delta(p-q)}{2},
\qquad
T_c=\frac{\delta(p+q+1)}{2},
\]
and the four transition vertices are exact cardinal half-steps:
\[
\begin{array}{c|cc}
&\Delta X&\Delta T\\
\hline
J_+&+\delta/2&0\\
J_-&-\delta/2&0\\
K_+&0&+\delta/2\\
K_-&0&-\delta/2.
\end{array}
\]

Consequently,
\begin{align*}
Y_{J+}^2
&=
T_c^2-\left(X_c+\frac{\delta}{2}\right)^2
=
\delta^2(p+1)q,\\
Y_{J-}^2
&=
T_c^2-\left(X_c-\frac{\delta}{2}\right)^2
=
\delta^2p(q+1),\\
Y_{K+}^2
&=
\left(T_c+\frac{\delta}{2}\right)^2-X_c^2
=
\delta^2(p+1)(q+1),\\
Y_{K-}^2
&=
\left(T_c-\frac{\delta}{2}\right)^2-X_c^2
=
\delta^2pq.
\end{align*}

Thus the same half-step
\[
\frac{\delta}{2}
\]
that completes the Casimir is literally the displacement from the state center to a transition edge in the side-view factor cell.  Squaring it produces
\[
\frac{\delta^2}{4}.
\]

\section{Relation to the unimodular null diamond}

The transition cell and the Farey null diamond are not identical objects.  Their exact common structure is the centered quadratic half-edge.

For a primitive arithmetic null edge,
\[
Q_\pm=M_F\pm D_F,
\]
with
\[
M_F^2=-D_F^2=\frac{\delta^2}{4}.
\]
For a ladder transition,
\[
Z_\pm=q\pm\frac{\delta}{2},
\]
and the Casimir completion is
\[
R^2=C+\frac{\delta^2}{4}.
\]

Both constructions therefore implement the same operation:
\[
\boxed{
\text{adjacent objects}
\longrightarrow
\text{edge center}
\longrightarrow
\left(\frac{\delta}{2}\right)^2.
}
\]

The arithmetic construction makes this invariant Lorentz-geometric.  The oscillator construction makes it representation-theoretic and dynamical.

\section{Cyclotomic normalization: a fourth quarter}

The discriminant-$12$ return supplies a fourth exact appearance of a quarter,
but in a different quadratic structure.  Let
\[
g_{12}=\begin{pmatrix}3&1\\2&1\end{pmatrix},
\qquad
H=g_{12}-\frac{\operatorname{tr}(g_{12})}{2}I
=\begin{pmatrix}1&1\\2&-1\end{pmatrix}.
\]
Then
\[
H^2=3I.
\]
After adjoining $i$, define the half-element
\begin{equation}
\mathcal Z=\frac{H+iI}{2}.
\label{eq:cyclotomic-half}
\end{equation}
Complex conjugation fixes $H$ and sends $i$ to $-i$, hence
\begin{equation}
\boxed{
\mathcal Z\overline{\mathcal Z}
=\frac{(H+iI)(H-iI)}{4}
=\frac{H^2+I}{4}
=I.
}
\label{eq:cyclotomic-quarter}
\end{equation}
The denominator $2$ is therefore not cosmetic: its square is exactly the
factor $4$ required to normalize the relative norm.  Equivalently, the two
eigenvalues of $\mathcal Z$ are
\[
\frac{\sqrt3+i}{2},\qquad \frac{-\sqrt3+i}{2},
\]
both primitive twelfth roots, and $\Phi_{12}(\mathcal Z)=0$.

At the integral level, if
\[
K=\mathbb Q(\sqrt3,i)=\mathbb Q(\zeta_{12}),
\qquad
\mathcal O_0=\mathbb Z[\sqrt3,i],
\]
then the maximal order satisfies
\[
\mathcal O_K/\mathcal O_0\cong(\mathbb F_2)^2.
\]
Thus cyclotomic closure adds a four-state parity gluing.  This statement and
\eqref{eq:cyclotomic-quarter} constitute the \emph{cyclotomic
normalization} of the quarter.  They do not identify the relative norm with
the Lorentz norm or with a representation-theoretic Casimir.

\begin{proposition}[Dyadic quarter taxonomy]
The following four exact constructions use the same elementary mechanism---a
canonical half-coordinate becomes a quarter under a quadratic operation:
\begin{enumerate}[label=(\roman*)]
\item the factor-cell midpoint square;
\item the compact/noncompact shifted Casimir;
\item the primitive null-edge midpoint norm;
\item the cyclotomic relative-norm normalization.
\end{enumerate}
Their common content is a mechanism type, not a universal operator identity.
\end{proposition}

\section{The ramified Boolean return}

The integral return also has an exact finite reduction.  Modulo $2$,
\[
\overline g_{12}
=\begin{pmatrix}1&1\\0&1\end{pmatrix},
\qquad
\overline g_{12}(x,y)=(x+y,y)
\quad\text{on }\mathbb F_2^2.
\]
This involution has one two-cycle and two fixed points.  It therefore cannot
be conjugate to either nonzero translation of the Boolean cell, since every
nonzero translation has two two-cycles.

The correct return is exposed by the Boolean relabeling
\[
\Phi(x,y)=(\varepsilon_1,\varepsilon_2)=(x,x+y).
\]
Indeed,
\begin{equation}
\boxed{
\Phi\,\overline g_{12}\,\Phi^{-1}
(\varepsilon_1,\varepsilon_2)
=(\varepsilon_2,\varepsilon_1).
}
\label{eq:boolean-exchange}
\end{equation}
Thus the ramified arithmetic return is factor exchange, not Boolean
translation.  If $\overline q_{12}(x,y)=y$ denotes the reduced degenerate norm,
then under the same relabeling
\begin{equation}
\boxed{
\overline q_{12}
=\varepsilon_1+\varepsilon_2
=\varepsilon_1\mathbin{\mathrm{xor}}\varepsilon_2.
}
\label{eq:boolean-xor}
\end{equation}

Let $V_4^{\mathrm{tr}}\cong(\mathbb F_2)^2$ be the translation group of the
four-state cell.  Adjoining the factor exchange in
\eqref{eq:boolean-exchange} gives
\[
\boxed{
V_4^{\mathrm{tr}}\rtimes C_2\cong D_8.
}
\]
This is the precise symmetry correction: the return lies in the affine
dihedral extension, not in the translation $V_4$ itself.

The finite space may be realized as the ramified quotient
\[
\mathfrak p_2/2\mathfrak p_2\cong(\mathbb F_2)^2.
\]
Although this quotient and $\mathcal O_K/\mathcal O_0$ both have four elements
and characteristic $2$, cardinality alone supplies no canonical isomorphism
between them.  Likewise, neither Boolean $V_4^{\mathrm{tr}}$ is to be confused
with the cyclotomic Galois group
$\operatorname{Gal}(\mathbb Q(\zeta_{12})/\mathbb Q)\cong V_4$.

\section{Scope of the interpretation}

Several distinctions are important.

First, the identities
\[
j(j+1)+\frac14=\left(j+\frac12\right)^2,
\qquad
k(k-1)+\frac14=\left(k-\frac12\right)^2
\]
are elementary completions of the square, and the appearance of $j+\tfrac12$ also has standard Weyl/Kirillov antecedents.  No claim is made that the Farey graph or multiplication-table geometry is the historical or representation-theoretic origin of that shift.

Second, the cone statement established here is stronger than a numerical coincidence.  The same $\delta^2/4$ is realized as:
\begin{enumerate}[label=(\roman*)]
\item the square of the half-spacing of adjacent weights;
\item the Lorentz norm squared of the midpoint of a scaled primitive unimodular null edge;
\item the quadratic displacement from the center of the two-mode transition cell to each cardinal transition edge.
\end{enumerate}
All three occur inside the same factor-cone quadratic form.

The cyclotomic quarter is an exact fourth realization, but it occurs in a
relative norm rather than in that factor-cone quadratic form.  Accordingly,
the common statement across all four cases is dyadic quadratic normalization,
not equality of the underlying quadratic spaces.

Third, the compact and noncompact ladder systems are not identified with one another.  Their distinction is preserved exactly by the sign in \eqref{eq:unified}, or geometrically by whether the completed Casimir radius occupies the timelike coordinate $T$ or the spacelike coordinate $X$.

\section{Conclusion}

The common quarter shift of the $SU(2)$ and $SU(1,1)$ quadratic Casimirs admits a simple geometric realization on the factor cone.  A primitive unimodular pair of null rays has an edge midpoint of squared Lorentz norm $1/4$, or $\delta^2/4$ at spacing $\delta$.  The exact rank-one ladder amplitudes use precisely this amount to pass from the quadratic Casimir to the edge-centered radius:
\[
\boxed{
R^2=C+M_F^2,
\qquad
M_F^2=\frac{\delta^2}{4}.
}
\]
The amplitudes then become cone sections:
\[
SU(2):\quad A^2=R^2-Z^2,
\qquad
SU(1,1):\quad A^2=Z^2-R^2.
\]

The resulting synthesis can be summarized as
\[
\boxed{
\text{Casimir completion}
\;\longleftrightarrow\;
\text{weight-edge centering}
\;\longleftrightarrow\;
\text{primitive null-edge centering}.
}
\]
The arithmetic null diamond should therefore be viewed not as the cause of the Casimir shift, but as a canonical arithmetic realization of the same universal edge-centered invariant.  The cyclotomic half-element supplies a distinct relative-norm completion,
\[
\mathcal Z\overline{\mathcal Z}=\frac{H^2+I}{4}=I,
\]
while the discriminant-$12$ return reduces to factor exchange with XOR norm on
a four-state ramified quotient.  Together these results give a four-part
taxonomy of the quarter while preserving the boundary between Casimir
completion, null-edge geometry, cyclotomic normalization, and Boolean return.

\section*{Source note}

The oscillator conventions used here are the standard Schwinger $SU(2)$ and two-mode $SU(1,1)$ realizations.  The author's current working quantum-realization manuscript remains under audit and is not used here as an independently audited theorem source.  Its factor-cone state dictionary is
\[
x=n_1+1,\qquad y=n_2,
\]
so that
\[
T=\frac{n_1+n_2+1}{2},\qquad
X=\frac{n_1-n_2+1}{2}.
\]
Within the four-corner transition-cell construction of this note, that state point is one transition vertex rather than the cell center.  For the two-mode positive discrete series, fixing the number difference $d=n_1-n_2$ gives the standard Bargmann index
\[
k=\frac{|d|+1}{2}.
\]
Only after choosing the oriented sector $d\ge0$ does this become
\[
k=\frac{d+1}{2}=X.
\]
The $SU(2)$ Casimir is $j(j+1)$ with $j=T-\tfrac12$.  The present note isolates and develops the edge-centered geometric consequence of these standard oscillator formulas as a standalone result, without depending on the unresolved broader Paper-C power-map claims.

\end{document}
